\section{提出假设}
\label{sec:hypotheses}

\subsection{理论依据}

中央银行对增长、通胀和金融风险的判断可能改变投资者对未来政策路径的预期。这些预期一方面进入折现率，另一方面影响对企业现金流和融资条件的判断。依赖外部融资的行业对信用条件更敏感，久期较长的行业可能对折现率更敏感。因此，同一份报告给市场提供的是共同信息，可能引起不同的行业反应。要检验行业差异，模型需要同时观察报告在时间上的措辞变化和行业在同一期的暴露差异。

报告语气可以同时反映经济判断和政策意向。当央行上调对经济形势的评价时，市场既可能上调盈利预期，也可能因预期政策支持减少而提高折现率。报告还会重述市场已经观察到的经济数据。在本文中，积极情绪与较高收益的关系是一项经验假设。词典按一致规则统计用词，适合提供可重复的基准。大语言模型能够利用句子语境解释否定、转折和限定语，其评分是否改善行业收益预测仍需与词典在相同样本上比较。

\subsection{研究假设与判定标准}

表~\ref{tab:hypothesis-rules} 给出本稿使用的判定标准。度量检验考察词典之间及词典与归档评分之间的一致性，事件检验考察发布窗口内的异常收益，季度模型则比较模型对行业收益排序和幅度的预测表现。信息系数（IC）为每季预测收益与实现收益的行业横截面秩相关。样本外决定系数（$R^2_{\mathrm{OOS}}$）和平均绝对误差（MAE）用于评价预测误差。下列阈值沿用大纲的暂定建议值，待团队与指导教师确认。本文据此给出初步判定。

\begin{table}[htbp]
  \centering\small
  \caption{研究假设及暂定判定标准}
  \label{tab:hypothesis-rules}
  \begin{tabularx}{\linewidth}{@{}l p{.22\linewidth} X l@{}}
    \toprule
    假设 & 待检验命题 & 判定依据 & 实验 \\
    \midrule
    H1 & 文本情绪可稳定度量 & 在 15 期归档评分中，主词典的 Pearson 或 Spearman 相关系数至少为 $0.50$，且未经校正的 $p<0.05$。主、次词典的 Pearson 相关系数至少为 $0.80$（103 期）。领域词典用于比较跨词表的一致性。同一报告无独立双评或间隔重测记录，评分者间一致性及重测稳定性均无法评估。 & E1 \\
    H2 & 更积极的措辞与更高的随后行业收益相关，并提供预测增量 & 事件维度考察主设定累计异常收益（CAR）及高、低情绪组差异，经多重校正后要求 $p<0.05$ 且情绪差异方向与假设一致（主样本 89 期）。预测维度要求共同测试样本上 IC 显著提高、$R^2_{\mathrm{OOS}}$ 提高且 MAE 下降（23 季）。 & E2--E4 \\
    H3 & 利率、融资或政策依赖程度不同的行业具有不同情绪暴露 & 联合模型中至少一个行业特征的差异项达到 $p<0.10$，且前后子样本符号一致率超过 $50\%$ 或秩相关显著（31 个行业）。 & E5 \\
    H4 & 大语言模型评分较词典有额外预测信息 & 共同测试样本的 IC 增幅达到 $p<0.05$（23 季），或联合模型中控制词典后，大语言模型的增量系数达到 $p<0.05$。 & E6 \\
    \bottomrule
  \end{tabularx}
\end{table}

\subsection{识别边界}

报告发布日期和措辞与当时的经济金融环境相关。事件窗内还可能有其他政策或宏观消息，市场模型只能剔除与基准收益线性相关的变动。本文将发布窗口的异常收益、跨期关联和季度排序表现分别报告。因果识别需要独立于市场同期预期的外生信息变动，本研究样本无法提供这一条件。

\section{实验及实验结果分析}
\label{sec:data}

\subsection{数据、变量与时间对齐}

表~\ref{tab:datasets} 汇总数据来源、覆盖范围及实证用途。中国人民银行网站提供《中国货币政策执行报告》（MPR）的详情页和 PDF 附件\cite{pbocMpr}。报告样本从 2001 年第一季度延续至 2026 年第三季度，共 103 期。索引同时记录报告期、发布日期、来源地址及下载文件的哈希值。行情取自 AkShare 转存的申万一级行业指数收盘价，覆盖 31 个行业，合计 144\,429 条日行情。按原始日行情的首条记录，16 个行业始于 1999 年 12 月 30 日，12 个始于 2014 年 2 月 21 日，石油石化、环保和美容护理始于 2021 年 12 月 13 日。早期事件的可用行业数因此较少。

\begin{table}[htbp]
  \centering\small
  \caption{实证数据及用途}
  \label{tab:datasets}
  \begin{tabularx}{\linewidth}{@{}p{.22\linewidth} X X@{}}
    \toprule
    数据 & 来源与覆盖 & 实证用途 \\
    \midrule
    货币政策执行报告 & 中国人民银行，2001Q1--2026Q3，103 期 & 发布时点、词典与大语言模型评分 \\
    申万一级行业指数 & AkShare 转存，1999-12-30 至 2026-11-17，31 个行业 & 日收益与季度收益、行业排序与组合 \\
    市场基准 & 沪深 300 和中证全指。E2/E3 使用的沪深 300 回溯序列始于 2002-01-04 & 交易日历、市场正常收益与基准替换 \\
    利率与宏观序列 & 中国 2 年期和 10 年期国债收益率、CPI、M2、PMI、社融等 & 超额市场模型、宏观控制与行业特征代理 \\
    中国三因子 & CH3 本地缓存，40 个覆盖完整的组合季度进入因子回归 & 多空收益的市场、规模和价值因子调整 \\
    \bottomrule
  \end{tabularx}
\end{table}

日简单收益定义为 $r_{i,t}=P_{i,t}/P_{i,t-1}-1$，其中 $P_{i,t}$ 为行业 $i$ 在交易日 $t$ 的收盘价。事件窗原始收益按日简单收益复利，累计异常收益按日异常收益求和。季度模型用各季最后可用收盘价构造相邻季的对数收益。事件日取官网发布日期，当日休市时顺延至下一可用交易日。若发布日期为交易日，现有日频设计将当日计入第 0 日，因此对于收盘后发布的报告，第 0 日收益可能来自公告前交易。

表~\ref{tab:alignment} 对齐报告发布日期与各收益窗口。本文将报告期下一自然季度的收益窗口简称为“原口径”，将报告发布后首个完整自然季度的收益窗口简称为“严格口径”。E4/E6 采用原口径。2026Q3 报告对应的 2026Q4 在行情截止时尚未完整，因此进入季度面板的是 2001Q2--2026Q3 的 102 个目标季。按这 102 期核对，发布日期均落入相应目标季，平均较季初晚 40.9 个自然日，跨度为 15 至 64 天。该模型的测试期在时间上留出，但目标收益包含发布前的价格变化。组合实验 E7 采用严格口径的持有窗口，沿用原口径模型给出的行业排序。第~\ref{sec:further} 章报告 E8 在严格口径下重新估计的结果。

\begin{table}[htbp]
  \centering\small
  \caption{报告发布日期与三类收益窗口的对齐}
  \label{tab:alignment}
  \begin{tabularx}{\linewidth}{@{}p{.25\linewidth} X@{}}
    \toprule
    分析口径 & 时间顺序与收益区间 \\
    \midrule
    事件窗（E2/E3） & 将发布日对齐至第 0 个交易日，用此前 120 个交易日估计正常收益，从第 0 日累计至第 1、3、5 或 10 日。 \\
    原口径（E4/E6） & 报告期结束 $\longrightarrow$ 目标季度开始 $\longrightarrow$ 报告发布 $\longrightarrow$ 目标季度结束。目标收益覆盖整个季度。 \\
    严格口径（E7/E8） & 报告发布 $\longrightarrow$ 下一自然季度开始 $\longrightarrow$ 该季度结束。E7 后移组合持有期，E8 按此窗口构造预测目标。 \\
    \bottomrule
  \end{tabularx}
\end{table}

\subsection{数据链路与词典度量（E0--E1）}

E0 以五期报告、三个行业和一个市场基准验证下载、文本提取、日历匹配与收益构造。五期均完成 PDF 获取、文本提取和发布日匹配，核验的 75 个行业收益单元全部可构造。参考交易日历中，三个行业于 2026 年 10 月 15 日各缺一条记录，总缺失率为 $3/1816=0.165\%$。后续计算保留缺失标记，事件研究按估计窗和事件窗的完整性筛选样本。E1 将文本管线扩展至 103 期，所有报告均取得 PDF 及可用文本，每期提取字符数超过 10\,000。

文本按行进行中文分词，过滤停用词、纯数字和标点。对 NTUSD 简体词典、大连理工大学情感词汇本体（DLUT）和自制领域词典分别剔除正负极性冲突词条。设 $N_{q,d}^{+}$ 和 $N_{q,d}^{-}$ 为报告 $q$ 在词典 $d$ 下的正向、负向词频，正文采用的归一化得分为
\begin{equation}
  S_{q,d}=\frac{N_{q,d}^{+}-N_{q,d}^{-}}{N_{q,d}^{+}+N_{q,d}^{-}}.
  \label{eq:dictionary-score}
\end{equation}
式~\eqref{eq:dictionary-score} 的分母为零时，得分记为零。NTUSD 简体词条等权，其词频得分与加权得分相同。脚本另保存 DLUT 和领域词典按各词表权重计算的得分。词典覆盖率以正、负极性词频之和除以清洗后的总词数，以衡量各词表参与评分的范围。

三套词典在全部报告中均有命中，平均覆盖率依次为 $10.79\%$、$6.78\%$ 和 $6.91\%$，同一文本在固定词表及分词规则下重复运行的分差为零。NTUSD 与 DLUT 的 Pearson 相关系数为 $0.829$（$p<0.001$），二者各自与自制词典的相关系数为 $0.349$ 和 $0.240$。图~\ref{fig:lexicon-series} 显示了主词典得分及覆盖率随时间的变动。三套词典对每期的得分均为正，因此符号一致率对区别词表的跨期排序能力帮助有限。未命中词表的汇总文件收录 17\,883 个不同词，其中包含“利率”“贷款”等依赖上下文判断极性的词。提高命中率本身未必提高情绪区分度。

\paperfigure{fig13_sentiment.png}{主词典 NTUSD 的报告情绪得分与覆盖率（E1）}{fig:lexicon-series}{0.88}

抽检以整篇报告为评分单位，采用从 $-2$ 到 $+2$ 的五级标尺评价整体语气。表~\ref{tab:human-agreement} 合并 5 期试评和 10 期正式抽检。归档记录标注三名评分人分别评了 6、5、4 期，每期仅有一份分值，且全部标记为已完成。归档表缺少评分日期、耗时及盲评记录，评分过程及独立性仍待核验，也缺少同一期报告的独立双评及间隔重测。E6 的临时评分参照另行使用。

NTUSD 和 DLUT 与归档分值的两种相关均为正，原始 $p$ 值均低于 $0.05$。领域词典的 Pearson 点估计达到 $0.50$，两种检验的原始 $p$ 值均高于 $0.05$。主词典达到表~\ref{tab:hypothesis-rules} 的暂定数值标准，H1 仅获初步、部分支持。人工基准的独立效度、评分者间一致性与重测稳定性仍待检验。

\begin{table}[htbp]
  \centering\small
  \caption[词典评分与归档分值的相关性]{词典评分与归档分值的相关性（15 期，每期单评分，评分独立性待核验）}
  \label{tab:human-agreement}
  \begin{tabular}{@{}lrrrr@{}}
    \toprule
    词典 & Pearson $r$ & $p$ 值 & Spearman $\rho$ & $p$ 值 \\
    \midrule
    NTUSD 简体 & 0.6954 & 0.0040 & 0.6334 & 0.0112 \\
    DLUT & 0.5995 & 0.0182 & 0.5756 & 0.0247 \\
    自制领域词典 & 0.5056 & 0.0545 & 0.4643 & 0.0812 \\
    \bottomrule
  \end{tabular}
  \par\smallskip
  \parbox{.92\linewidth}{\footnotesize 注：本表合并试评和正式抽检。每期只有一份归档分值，评分过程及独立性待核验。三个词典的相关检验均未做多重校正。}
\end{table}

\subsection{事件研究（E2--E3）}

参照市场模型的事件研究设计\cite{macKinlay1997}，在每次事件的 $[-120,-1]$ 交易日内，为各行业回归估计正常收益 $\widehat r_{i,t}=\widehat\alpha_i+\widehat\beta_i r_{m,t}$，其中 $r_{m,t}$ 为沪深 300 日收益率。设事件于交易日 $t_0$ 对齐，行业 $i$ 的累计异常收益为 $\mathrm{CAR}_{i}([0,k])=\sum_{j=0}^{k}(r_{i,t_0+j}-\widehat r_{i,t_0+j})$。对照模型分别使用估计窗内行业日收益均值和扣除 2 年期国债日利率的超额市场收益。

先在每期对可用行业的 CAR 等权平均，再以报告期为单位做推断。NTUSD 以全部报告分数的秩三等分形成高、中、低组，组别表示相对语气。与零比较使用单样本 $t$ 检验及符号检验，组间差异使用 Welch $t$ 检验，置信区间用 10\,000 次期层面自助抽样计算。103 期中 95 期发布日为交易日、5 期顺延、3 期早于主基准起点。应用窗长和行情完整性要求后，主模型有 89 个可用事件，每期平均约 19.4 个行业。E3 日志跨模型和窗长登记 5\,851 条跳过记录，计数包含不同层级的缺失或不完整记录。

\begin{table}[htbp]
  \centering\small
  \caption{市场模型事件窗的累计异常收益（E3）}
  \label{tab:event-results}
  \setlength{\tabcolsep}{4pt}
  \begin{tabular}{@{}crrrrrrrr@{}}
    \toprule
    & \multicolumn{4}{c}{全样本（89 期）} & \multicolumn{4}{c}{高减低（29 对 30 期）} \\
    \cmidrule(lr){2-5}\cmidrule(lr){6-9}
    事件窗 & CAR（\%） & 原始 $p$ & BH $p$ & Holm $p$ & 差值（百分点） & 原始 $p$ & BH $p$ & Holm $p$ \\
    \midrule
    $[0,+1]$  & 0.0827 & 0.4831 & 0.6408 & 1.0000 & $-0.6089$ & 0.0113 & 0.0906 & 0.0906 \\
    $[0,+3]$  & 0.2896 & 0.1479 & 0.2957 & 0.7393 & $-0.2615$ & 0.5751 & 0.6408 & 1.0000 \\
    $[0,+5]$  & 0.3999 & 0.0696 & 0.1855 & 0.4173 & $-0.5837$ & 0.2503 & 0.4005 & 1.0000 \\
    $[0,+10]$ & 0.6530 & 0.0409 & 0.1637 & 0.2865 & $-0.3438$ & 0.6408 & 0.6408 & 1.0000 \\
    \bottomrule
  \end{tabular}
  \par\smallskip
  \parbox{.92\linewidth}{\footnotesize 注：用 E3 原始输出中的四个全样本和四个高减低检验构成一个检验族。Benjamini--Hochberg（BH）与 Holm 校正后的最小 $p$ 值均为 $0.0906$。E8 按其设定另行划分稳健性检验族。}
\end{table}

图~\ref{fig:event-car} 给出各情绪组的 CAR 及置信区间。全样本 $[0,+10]$ 窗平均 CAR 为 $0.653\%$（原始 $p=0.0409$）。相同窗口的超额市场模型原始 $p=0.0235$（81 期），均值调整模型各全样本窗口均未达 $5\%$ 水平。低情绪组 $[0,+1]$ 窗平均 CAR 为 $0.456\%$（原始 $p=0.0219$），而高减低差值为 $-0.609$ 个百分点（原始 $p=0.0113$），方向与 H2 预期相反。表~\ref{tab:event-results} 的检验经两种多重比较校正后均未达 $5\%$ 水平。本期估计窗与本期事件窗不重叠。在 120 日设定的重叠核查中，可用估计窗均覆盖上一期 MPR 发布日。最大事件窗内另有 17 期被标记为同期冲击，剔除后全样本 $[0,+10]$ 窗原始 $p$ 上升至 $0.0702$。这些结果提示短窗反应依赖窗口和样本选择，H2 的事件维度未获支持。

\paperfigure{fig14_event_study_car.png}{市场模型下各情绪组的累计异常收益（E3）。阴影为未经多重校正的 95\% 自助置信区间，本期估计窗与本期事件窗不重叠}{fig:event-car}{0.82}

\subsection{行业收益的描述性统计}

行业间的风险收益比较以 2014 年 2 月 21 日至 2026 年 11 月 17 日为共同日历区间，排除在 2021 年才开始记录的三个行业。先在各行业完整价格序列上计算日收益，再按日期筛选。以 252 个交易日换算年化波动率和复利年化收益，最大回撤则从累计净值的历史峰值计算。图~\ref{fig:risk-return} 展示 28 个可比行业的年化收益与年化波动率。复利年化收益介于 $-5.41\%$（商贸零售）与 $10.12\%$（煤炭），年化波动率介于 $21.44\%$（银行）与 $35.28\%$（煤炭）。银行的最大回撤为 $-37.44\%$，商贸零售为 $-84.33\%$。行业日收益的样本超额峰度介于 2.17 和 7.02。收益及风险的离散程度提示季度排序需要控制行业特征，并以共同观测比较模型。

\paperfigure{fig09_risk_return.png}{28 个可比行业的年化收益与年化波动率（2014-02-21 至 2026-11-17）}{fig:risk-return}{0.78}

E1 的结果为季度情绪变量提供了初步度量依据。E2/E3 的结果则表明，个别未校正的异常收益不足以识别稳定的情绪方向。下一章用行业暴露交互项检验共同文本信息能否改善季度收益排序。
